\chapter{Complexity — Universal Startup}
%%% generated by the blueprint pipeline from TCSlib.Complexity.TuringMachine.UniversalStartup
%%%%%%%%%%%%%%%%%%%%%%%%%%%%%%%%%%%%%%%%%%%%%%%%%%%%%%%%%%%%%%%%%%%%%%%%%%%%%%%
\section{Overview}

This module builds the startup layer of the universal Turing machine
(Arora--Barak, \S 1.4.1, Theorem 1.9): it parses the doubled-bit code region
of the self-delimiting pairing $\langle \alpha, x\rangle$ of a program code
$\alpha$ and an input $x$, runs the coding scheme's canonizer on the extracted
code, embeds a virtual input head into the verbatim suffix region, and
captures the resulting canonical transition table on a dedicated work tape
before control transfers to the table interpreter.  Throughout, bit strings
are written over $\{0,1\}$ ($1$ for true, $0$ for false), and
$\langle \alpha, x\rangle$ lists each bit of $\alpha$ twice, then the
separator pair $01$, then $x$ verbatim.

%%%%%%%%%%%%%%%%%%%%%%%%%%%%%%%%%%%%%%%%%%%%%%%%%%%%%%%%%%%%%%%%%%%%%%%%%%%%%%%
\section{Declarations}

\begin{lemma}[Length of the paired encoding]
\lean{Turing.universal_pair_length}
\label{Turing.universal_pair_length}
\leanok
\uses{Turing.pairEncode}
\difficulty{3}
For all bit strings $\alpha$ and $x$, the self-delimiting pairing
$\langle \alpha, x\rangle$ --- each bit of $\alpha$ written twice, followed by
the separator pair $01$, followed by $x$ verbatim --- has length exactly
$2\abs{\alpha} + 2 + \abs{x}$.  In particular the extent of the code region is
independent of the suffix $x$.
\end{lemma}

\begin{lemma}[Aligned cells of a doubled code bit]
\lean{Turing.universal_pair_get}
\label{Turing.universal_pair_get}
\leanok
\uses{Turing.pairEncode, Turing.universalPrefixTM}
\difficulty{4}
Let $\alpha$ and $x$ be bit strings and let $j < \abs{\alpha}$.  In the
self-delimiting pairing $\langle \alpha, x\rangle$ (each bit of $\alpha$
doubled, then the separator pair $01$, then $x$ verbatim), the cells of index
$2j$ and $2j+1$ (zero-indexed) both contain the $j$-th bit of $\alpha$; both
cells of a doubled code bit therefore lie before the separator.
\end{lemma}

\begin{lemma}[Location of the separator]
\lean{Turing.universal_pair_separator}
\label{Turing.universal_pair_separator}
\leanok
\uses{Turing.pairEncode, Turing.universalPrefixTM}
\difficulty{3}
For all bit strings $\alpha$ and $x$, the cells of index $2\abs{\alpha}$ and
$2\abs{\alpha}+1$ in the self-delimiting pairing $\langle \alpha, x\rangle$
(each bit of $\alpha$ doubled, then the separator pair $01$, then $x$
verbatim) contain the bits $0$ and $1$ respectively: the separator begins
immediately after the doubled code region.
\end{lemma}

\begin{definition}[The doubled-bit prefix parser]
\lean{Turing.universalPrefixTM}
\label{Turing.universalPrefixTM}
\leanok
\uses{Turing.FinTM}
A Turing machine over the binary alphabet with no work tapes and two kinds of
live control states: ``expecting the first half of a doubled pair'' and
``having just read a first half $b$''.  In the expecting state it records the
scanned bit and moves right (halting if it scans a blank); in the state
remembering $b$ it compares the scanned cell with $b$ --- on a match it emits
the single bit $b$, moves right, and returns to the expecting state, while on
a mismatch (which on a valid paired encoding occurs exactly at the
separator's second bit) it moves right and halts without emitting.
\end{definition}

\begin{definition}[Configurations of the prefix parser]
\lean{Turing.universalPrefixCfg}
\label{Turing.universalPrefixCfg}
\leanok
\uses{Turing.Cfg, Turing.pairEncode}
Given bit strings $\alpha$ and $x$, a control state, an input-head position on
the self-delimiting pairing $\langle \alpha, x\rangle$, and an output word,
this assembles the corresponding full configuration of the doubled-bit prefix
parser; since that parser has no work tapes, the work-tape and work-head
components are vacuous.
\end{definition}

\begin{lemma}[Single transition of the prefix parser]
\lean{Turing.universalPrefix_step}
\label{Turing.universalPrefix_step}
\leanok
\uses{Turing.Action.apply, Turing.Cfg.ext_zero_tapes, Turing.Cfg.inputSymbol, Turing.Cfg.workTapeSymbols, Turing.MultiTapeTM.step, Turing.moveInputPos, Turing.pairEncode, Turing.universalPrefixCfg, Turing.universalPrefixTM}
\difficulty{3}
Let the doubled-bit prefix parser stand, on the self-delimiting pairing
$\langle \alpha, x\rangle$ of bit strings $\alpha$ and $x$, in a live control
state $q$ at input-head position $p$ with output word $w$ emitted so far, and
suppose the symbol scanned at $p$ is $b$.  Then one transition produces again
a configuration with no work-tape content: its control state and head
movement are those prescribed by the parser's transition function on
$(q, b)$, and the bit emitted by that transition (if any) is appended to $w$.
\end{lemma}

\begin{lemma}[Prefix parser after an even number of steps]
\lean{Turing.universalPrefix_bits}
\label{Turing.universalPrefix_bits}
\leanok
\uses{Turing.Cfg.ext_zero_tapes, Turing.FinTM.inputSymbol_at, Turing.MultiTapeTM.initCfg, Turing.MultiTapeTM.runFrom, Turing.MultiTapeTM.runFrom_succ_eq_step', Turing.moveInputPos_pos_of_ne_right, Turing.pairEncode, Turing.universalPrefixCfg, Turing.universalPrefixTM, Turing.universalPrefix_step, Turing.universal_pair_get, Turing.universal_pair_length}
\difficulty{5}
Let $\alpha$ and $x$ be bit strings and let $j \le \abs{\alpha}$.  Started on
the self-delimiting pairing $\langle \alpha, x\rangle$ (each bit of $\alpha$
doubled, then the separator pair $01$, then $x$ verbatim), the doubled-bit
prefix parser is, after exactly $2j$ transitions, in the ``expecting a first
half'' state with its input head at position $2j+1$ --- head positions count
the left boundary blank as position $0$, so the head stands on the cell of
index $2j$ --- having emitted precisely the first $j$ bits of $\alpha$.
\end{lemma}

\begin{lemma}[Completion of prefix extraction]
\lean{Turing.universalPrefix_start}
\label{Turing.universalPrefix_start}
\leanok
\uses{Turing.Cfg.ext_zero_tapes, Turing.FinTM.inputSymbol_at, Turing.MultiTapeTM.initCfg, Turing.MultiTapeTM.runFrom, Turing.MultiTapeTM.runFrom_succ_eq_step', Turing.moveInputPos_pos_of_ne_right, Turing.pairEncode, Turing.universalPrefixCfg, Turing.universalPrefixTM, Turing.universalPrefix_bits, Turing.universalPrefix_step, Turing.universal_pair_length, Turing.universal_pair_separator}
\difficulty{4}
For all bit strings $\alpha$ and $x$, the doubled-bit prefix parser, started
on the self-delimiting pairing $\langle \alpha, x\rangle$, halts after
exactly $2\abs{\alpha}+2$ transitions with output exactly $\alpha$ and with
its input head parked at position $2\abs{\alpha}+3$, that is, on the first
cell of the verbatim suffix $x$ (the right blank when $x$ is empty).  Both
the step count and the emitted word are independent of $x$.
\end{lemma}

\begin{lemma}[Prefix parser one step before halting]
\lean{Turing.universalPrefix_penultimate}
\label{Turing.universalPrefix_penultimate}
\leanok
\uses{Turing.FinTM.inputSymbol_at, Turing.MultiTapeTM.initCfg, Turing.MultiTapeTM.runFrom, Turing.MultiTapeTM.runFrom_succ_eq_step', Turing.moveInputPos_pos_of_ne_right, Turing.pairEncode, Turing.universalPrefixCfg, Turing.universalPrefixTM, Turing.universalPrefix_bits, Turing.universalPrefix_step, Turing.universal_pair_length, Turing.universal_pair_separator}
\difficulty{3}
For all bit strings $\alpha$ and $x$, after exactly $2\abs{\alpha}+1$
transitions on the self-delimiting pairing $\langle \alpha, x\rangle$, the
doubled-bit prefix parser is in the live state recording the separator's
first bit $0$, its input head is at position $2\abs{\alpha}+2$, and it has
emitted exactly $\alpha$.
\end{lemma}

\begin{lemma}[Liveness propagates backwards]
\lean{Turing.universal_live_before}
\label{Turing.universal_live_before}
\leanok
\uses{Turing.Cfg, Turing.MultiTapeTM, Turing.MultiTapeTM.runFrom, Turing.MultiTapeTM.runFrom_add, Turing.MultiTapeTM.runFrom_of_halt}
\difficulty{3}
Let $M$ be a multi-tape Turing machine, $c$ one of its configurations on some
input, and $s \le t$ natural numbers.  If the configuration reached from $c$
after $t$ steps is not halted, then neither is the configuration reached
after $s$ steps: since halting is absorbing, a live later configuration
excludes all earlier halts.
\end{lemma}

\begin{lemma}[Liveness of prefix extraction]
\lean{Turing.universalPrefix_live}
\label{Turing.universalPrefix_live}
\leanok
\uses{Turing.MultiTapeTM.initCfg, Turing.MultiTapeTM.runFrom, Turing.pairEncode, Turing.universalPrefixTM, Turing.universalPrefix_penultimate, Turing.universal_live_before}
\difficulty{3}
For all bit strings $\alpha$ and $x$ and every $s < 2\abs{\alpha}+2$, the
doubled-bit prefix parser, started on the self-delimiting pairing
$\langle \alpha, x\rangle$, has not halted after $s$ transitions.
\end{lemma}

\begin{definition}[Prefix extraction composed with the canonizer]
\lean{Turing.universalCanonTM}
\label{Turing.universalCanonTM}
\leanok
\uses{Turing.EffectiveMachineCode, Turing.FinTM, Turing.FinTM.bufferedCompTM, Turing.universalPrefixTM}
Given an effective machine-coding scheme --- a representation of machines by
bit strings equipped with a canonizer machine that computes, from any code,
the serialization of the canonical transition table of the decoded machine
within a stated time bound --- this is the buffered composition of the
doubled-bit prefix parser with that canonizer: the word emitted by the parser
(the extracted code) is stored in a buffer, and the canonizer then runs with
its input served from that buffer while the physical input head stays put.
\end{definition}

\begin{lemma}[Canonizer start after prefix extraction]
\lean{Turing.universalCanon_start}
\label{Turing.universalCanon_start}
\leanok
\uses{Turing.EffectiveMachineCode, Turing.FinTM.bufferedCompTM, Turing.FinTM.bufferedFirstCfg_init, Turing.FinTM.bufferedFirstCfg_rewind, Turing.FinTM.bufferedFirstCfg_run, Turing.FinTM.bufferedSecondCfg, Turing.MultiTapeTM.initCfg, Turing.MultiTapeTM.runFrom, Turing.MultiTapeTM.runFrom_add, Turing.pairEncode, Turing.universalCanonTM, Turing.universalPrefixTM, Turing.universalPrefix_live, Turing.universalPrefix_start, Turing.universal_pair_length}
\difficulty{4}
Let $c$ be an effective machine-coding scheme and let $\alpha$ and $x$ be bit
strings.  Started on the self-delimiting pairing $\langle \alpha, x\rangle$,
the buffered composition of the doubled-bit prefix parser with the canonizer
of $c$ reaches, after exactly $3\abs{\alpha}+4$ transitions, the second-phase
buffered configuration in which the canonizer stands at its initial
configuration on input $\alpha$, the virtual-boundary flag is set, and the
physical input head is parked at position $2\abs{\alpha}+3$, on the first
cell of the verbatim suffix.  This step count does not depend on $x$.
\end{lemma}

\begin{lemma}[The canonizer run is served from the buffer]
\lean{Turing.universalCanon_run}
\label{Turing.universalCanon_run}
\leanok
\uses{Turing.EffectiveMachineCode, Turing.FinTM.VirtualTag, Turing.FinTM.bufferedSecondCfg, Turing.FinTM.bufferedSecondCfg_run, Turing.MultiTapeTM.initCfg, Turing.MultiTapeTM.runFrom, Turing.MultiTapeTM.runFrom_add, Turing.pairEncode, Turing.universalCanonTM, Turing.universalCanon_start, Turing.universalPrefixTM, Turing.universal_pair_length}
\difficulty{4}
Let $c$ be an effective machine-coding scheme, $\alpha$ and $x$ bit strings,
and $t$ a natural number.  Then there is a Boolean boundary flag $b$ such
that, after $3\abs{\alpha}+4+t$ transitions on the self-delimiting pairing
$\langle \alpha, x\rangle$, the buffered composition of the doubled-bit
prefix parser with the canonizer of $c$ is in the second-phase buffered
configuration embedding the configuration reached by the canonizer after $t$
steps on input $\alpha$, with flag $b$ and with the physical input head still
parked at position $2\abs{\alpha}+3$.
\end{lemma}

\begin{lemma}[Canonical table at a suffix-independent time]
\lean{Turing.universalCanon_complete}
\label{Turing.universalCanon_complete}
\leanok
\uses{Turing.CodeTM.serialize, Turing.EffectiveMachineCode, Turing.FinTM.bufferedSecondCfg, Turing.FinTM.computesInTime_iff, Turing.MultiTapeTM.initCfg, Turing.MultiTapeTM.runFrom, Turing.pairEncode, Turing.universalCanonTM, Turing.universalCanon_run}
\difficulty{3}
Let $c$ be an effective machine-coding scheme with canonizer time bound
$T_c$, and let $\alpha$ and $x$ be bit strings.  After exactly
$3\abs{\alpha}+4+T_c(\abs{\alpha})$ transitions on the self-delimiting
pairing $\langle \alpha, x\rangle$, the buffered composition of the
doubled-bit prefix parser with the canonizer of $c$ has halted, its output is
the serialization of the canonical transition table of the machine that $c$
decodes from $\alpha$, and its input head is at position $2\abs{\alpha}+3$.
In particular this time depends only on the code, not on the suffix $x$.
\end{lemma}

\begin{lemma}[Suffix lookups in the paired encoding]
\lean{Turing.universal_pair_suffix}
\label{Turing.universal_pair_suffix}
\leanok
\uses{Turing.pairEncode, Turing.universal_pair_length}
\difficulty{3}
For all bit strings $\alpha$ and $x$ and every natural number $j$, the
optional cell lookup at index $2\abs{\alpha}+2+j$ in the self-delimiting
pairing $\langle \alpha, x\rangle$ agrees with the optional lookup at index
$j$ in $x$; both are simultaneously undefined when $j \ge \abs{x}$.
\end{lemma}

\begin{definition}[Embedding of the virtual suffix head]
\lean{Turing.universalInputPos}
\label{Turing.universalInputPos}
\leanok
\uses{Turing.pairEncode, Turing.universal_pair_length}
Given bit strings $\alpha$ and $x$, this maps a virtual input-head position
$p \in \{0, 1, \dots, \abs{x}+1\}$ on the suffix $x$ (including both boundary
blanks) to the physical head position $2\abs{\alpha}+2+p$ on the
self-delimiting pairing $\langle \alpha, x\rangle$.  Virtual position $0$
lands on the separator's last cell.
\end{definition}

\begin{definition}[Configuration embedding along the suffix]
\lean{Turing.universalInputCfg}
\label{Turing.universalInputCfg}
\leanok
\uses{Turing.Cfg, Turing.pairEncode, Turing.universalInputPos}
Given a bit string $\alpha$ and a configuration of a machine with $k$ work
tapes on input $x$, this forms the configuration on the self-delimiting
pairing $\langle \alpha, x\rangle$ that keeps the control state, work tapes,
work-head positions, and output unchanged and translates only the input-head
position by the offset $2\abs{\alpha}+2$ of the verbatim suffix.
\end{definition}

\begin{lemma}[Reading the left-boundary marker tape]
\lean{Turing.universal_marker}
\label{Turing.universal_marker}
\leanok
\uses{Turing.FinTM.bufferTape, Turing.FinTM.bufferTape_nat}
\difficulty{2}
Consider the tape that carries the single-cell word $1$ at coordinate $0$
and is blank at every other coordinate, including all negative ones.  For
every natural number $p$, the cell of this tape at coordinate $p$ contains
the mark $1$ if $p = 0$ and is blank otherwise.
\end{lemma}

\begin{lemma}[Masked read equals the native input symbol]
\lean{Turing.universalInput_read}
\label{Turing.universalInput_read}
\leanok
\uses{Turing.Cfg, Turing.Cfg.inputSymbol, Turing.FinTM.bufferTape, Turing.FinTM.inputSymbol_at, Turing.universalInputCfg, Turing.universalInputPos, Turing.universal_marker, Turing.universal_pair_length, Turing.universal_pair_suffix}
\difficulty{4}
Let $\alpha$ be a bit string and let $c$ be a configuration of a machine with
$k$ work tapes on input $x$; embed $c$ into the self-delimiting pairing
$\langle \alpha, x\rangle$ by shifting its input head by the suffix offset
$2\abs{\alpha}+2$.  Mask the embedded read as follows: if the marker tape
carrying a single mark at coordinate $0$ shows its mark at the coordinate
equal to $c$'s virtual head position --- that is, exactly at virtual position
$0$ --- declare the read blank; otherwise take the symbol physically scanned
by the embedded configuration.  This masked read equals the input symbol that
$c$ scans natively on $x$, including at the left and right boundary blanks
(and for an empty suffix).
\end{lemma}

\begin{lemma}[The marker yields a valid boundary tag]
\lean{Turing.universalInput_tag}
\label{Turing.universalInput_tag}
\leanok
\uses{Turing.FinTM.VirtualTag, Turing.FinTM.bufferTape, Turing.universal_marker}
\difficulty{2}
For every head position $p \in \{0, 1, \dots, n+1\}$ on an input of length
$n$ (with its two boundary blanks), the Boolean flag ``the marker tape
carrying a single mark at coordinate $0$ does not show its mark at coordinate
$p$'' satisfies the boundary-tag contract of the clamped virtual-move
operation: the flag is false at the left boundary $p = 0$ and true at the
right blank $p = n+1$, including when the input is empty.
\end{lemma}

\begin{lemma}[Virtual motion commutes with the suffix embedding]
\lean{Turing.universalInput_move}
\label{Turing.universalInput_move}
\leanok
\uses{Turing.Cfg, Turing.Cfg.inputSymbol, Turing.FinTM.bufferTape, Turing.FinTM.virtualMove, Turing.FinTM.virtualMove_correct, Turing.moveInputPos, Turing.universalInputPos, Turing.universalInput_tag, Turing.universal_pair_length}
\difficulty{5}
Let $\alpha$ be a bit string, let $c$ be a configuration of a machine with
$k$ work tapes on input $x$ whose input head is at position $p$, and let
$d \in \{-1, 0, +1\}$ be a direction.  Let $b$ be the marker-derived boundary
flag at $p$ (false exactly at $p = 0$) and let $m$ be the clamped virtual
move determined by $b$, the symbol under $c$'s input head, and $d$, which
suppresses stepping outward from a boundary blank.  Then moving the physical
head from the embedded position $2\abs{\alpha}+2+p$ on the self-delimiting
pairing $\langle \alpha, x\rangle$ by $m$ lands exactly on the embedding of
the position reached from $p$ by the clamped move in direction $d$, and, as
integers, $p + m$ equals that new virtual position.
\end{lemma}

\begin{definition}[The output-capture wrapper]
\lean{Turing.universalCaptureTM}
\label{Turing.universalCaptureTM}
\leanok
\uses{Turing.FinTM, Turing.FinTM.controlAction, Turing.FinTM.rightAction, Turing.FinTM.tapeBlocks, Turing.MultiTapeTM}
Given a machine $M$ over the binary alphabet with $k$ work tapes and a
four-tape machine $D$ (the interpreter), the wrapper machine with $k+4$ work
tapes and two-phase control.  In the first phase it runs $M$ on the first $k$
work tapes and the real input, except that every bit $M$ emits is written
onto a dedicated table tape --- the first of the four extra tapes --- instead
of the real output.  When $M$ halts, a single administrative transition
passes control to $D$, which then runs on the four extra tapes, reads the
original physical input, and emits output for real; the three extra tapes
other than the table tape start blank.
\end{definition}

\begin{definition}[First-phase configurations of the capture wrapper]
\lean{Turing.universalCaptureCfg}
\label{Turing.universalCaptureCfg}
\leanok
\uses{Turing.Cfg, Turing.FinTM, Turing.FinTM.bufferTape, Turing.FinTM.tapeBlocks, Turing.MultiTapeTM, Turing.universalCaptureTM}
For a configuration $c$ of a machine $M$ with $k$ work tapes on input $x$,
the corresponding first-phase configuration of the output-capture wrapper of
$M$ and an interpreter $D$: the control records phase one together with $c$'s
control state, the input head and the first $k$ work tapes are those of $c$,
the table tape holds the word output by $c$ so far (written from coordinate
$0$) with its head one cell past the end of that word, the remaining three
tapes are blank with heads at $0$, and the wrapper's real output is empty.
\end{definition}

\begin{lemma}[Initialization of the capture wrapper]
\lean{Turing.universalCapture_init}
\label{Turing.universalCapture_init}
\leanok
\uses{Turing.FinTM, Turing.FinTM.tapeBlocks, Turing.MultiTapeTM, Turing.MultiTapeTM.initCfg, Turing.universalCaptureCfg, Turing.universalCaptureTM}
\difficulty{4}
For every machine $M$, four-tape interpreter $D$, and input $x$, the initial
configuration of the output-capture wrapper of $M$ and $D$ on $x$ coincides
with the first-phase embedding of the initial configuration of $M$ on $x$:
blank table tape, blank interpreter tapes, and empty output.
\end{lemma}

\begin{lemma}[One captured transition]
\lean{Turing.universalCapture_step}
\label{Turing.universalCapture_step}
\leanok
\uses{Turing.Action, Turing.Action.apply, Turing.Cfg, Turing.Cfg.inputSymbol, Turing.Cfg.workTapeSymbols, Turing.FinTM, Turing.FinTM.bufferTape_append, Turing.FinTM.tapeBlocks, Turing.MultiTapeTM, Turing.MultiTapeTM.step, Turing.universalCaptureCfg, Turing.universalCaptureTM, Turing.universalCapture_run}
\difficulty{5}
Let $c$ be a configuration of a machine $M$ on input $x$ whose control state
is live (not halted).  Then one transition of the output-capture wrapper of
$M$ and a four-tape interpreter $D$, applied to the first-phase embedding of
$c$, equals the first-phase embedding of the configuration obtained from $c$
by one transition of $M$.  In particular any bit emitted by that transition
of $M$ --- including a bit emitted on $M$'s halting transition --- is
appended on the table tape, and the wrapper itself remains live.
\end{lemma}

\begin{lemma}[Lockstep capture of a live run]
\lean{Turing.universalCapture_run}
\label{Turing.universalCapture_run}
\leanok
\uses{Turing.Cfg, Turing.FinTM, Turing.MultiTapeTM, Turing.MultiTapeTM.runFrom, Turing.MultiTapeTM.runFrom_succ_eq_step', Turing.universalCaptureCfg, Turing.universalCaptureTM, Turing.universalCapture_step}
\difficulty{3}
Let $c$ be a configuration of a machine $M$ on input $x$ and let $t$ be a
natural number such that $M$ has not halted after any $s < t$ steps from $c$.
Then $t$ transitions of the output-capture wrapper of $M$ and a four-tape
interpreter $D$, started from the first-phase embedding of $c$, produce
exactly the first-phase embedding of the configuration of $M$ after $t$
steps from $c$.
\end{lemma}

\begin{definition}[Interpreter-entry configuration]
\lean{Turing.universalCapturedCfg}
\label{Turing.universalCapturedCfg}
\leanok
\uses{Turing.Cfg, Turing.FinTM, Turing.MultiTapeTM, Turing.universalCaptureCfg, Turing.universalCaptureTM}
For a configuration $c$ of a machine $M$ on input $x$, the configuration of
the output-capture wrapper of $M$ and a four-tape interpreter $D$ that agrees
with the first-phase embedding of $c$ in every component except that its
control state is the interpreter's initial state.  The halted source
configuration is thereby preserved as inactive data, with its output stored
on the table tape.
\end{definition}

\begin{lemma}[Transfer to the interpreter]
\lean{Turing.universalCapture_transfer}
\label{Turing.universalCapture_transfer}
\leanok
\uses{Turing.Cfg, Turing.FinTM, Turing.MultiTapeTM, Turing.MultiTapeTM.step, Turing.moveInputPos_zero, Turing.universalCaptureCfg, Turing.universalCaptureTM, Turing.universalCapturedCfg}
\difficulty{3}
Let $c$ be a halted configuration of a machine $M$ on input $x$.  Then one
transition of the output-capture wrapper of $M$ and a four-tape interpreter
$D$, applied to the first-phase embedding of $c$, yields the
interpreter-entry configuration for $c$: every component of $c$ is preserved
and the control becomes the interpreter's initial state.
\end{lemma}

\begin{lemma}[Every completed run reaches the interpreter]
\lean{Turing.universalCapture_start}
\label{Turing.universalCapture_start}
\leanok
\uses{Turing.FinTM, Turing.MultiTapeTM, Turing.MultiTapeTM.initCfg, Turing.MultiTapeTM.runFrom, Turing.MultiTapeTM.runFrom_add, Turing.MultiTapeTM.runFrom_of_halt, Turing.MultiTapeTM.runFrom_succ_eq_step', Turing.universalCaptureTM, Turing.universalCapture_init, Turing.universalCapture_run, Turing.universalCapture_transfer, Turing.universalCapturedCfg}
\difficulty{5}
Let $M$ be a machine, $D$ a four-tape interpreter, $x$ an input, and $T$ a
natural number such that $M$ has halted within $T$ steps on $x$.  Then there
exists $t \le T+1$ such that, after exactly $t$ transitions on $x$, the
output-capture wrapper of $M$ and $D$ stands in the interpreter-entry
configuration corresponding to the configuration of $M$ after $T$ steps: the
control is the interpreter's initial state, the full halted configuration of
$M$ --- including its parked input head --- is preserved, and the output of
$M$ is captured on the table tape.
\end{lemma}

\begin{lemma}[Startup with the canonical table captured]
\lean{Turing.universal_captured_table}
\label{Turing.universal_captured_table}
\leanok
\uses{Turing.CodeTM.serialize, Turing.EffectiveMachineCode, Turing.FinTM.bufferTape, Turing.FinTM.tapeBlocks, Turing.FinTM.tapeBlocks_buffer, Turing.MultiTapeTM, Turing.MultiTapeTM.initCfg, Turing.MultiTapeTM.runFrom, Turing.pairEncode, Turing.universalCanonTM, Turing.universalCanon_complete, Turing.universalCaptureTM, Turing.universalCapture_start}
\difficulty{4}
Let $c$ be an effective machine-coding scheme with canonizer time bound
$T_c$, let $D$ be a four-tape interpreter, and let $\alpha$ and $x$ be bit
strings.  Start the output-capture wrapper of the composed
prefix-extraction/canonizer machine of $c$, with interpreter $D$, on the
self-delimiting pairing $\langle \alpha, x\rangle$.  Then there is a
$t \le 3\abs{\alpha}+5+T_c(\abs{\alpha})$ such that after $t$ transitions the
control is the interpreter's initial state, the physical input head is at
position $2\abs{\alpha}+3$ (the start of the verbatim suffix), the
interpreter's first work tape holds --- written from coordinate $0$ and blank
elsewhere --- the serialization of the canonical transition table of the
machine that $c$ decodes from $\alpha$, and the real output is still empty.
\end{lemma}
